%% bare_jrnl.tex
%% V1.4b
%% 2015/08/26
%% by Michael Shell
%% see http://www.michaelshell.org/
%% for current contact information.
%%
%% This is a skeleton file demonstrating the use of IEEEtran.cls
%% (requires IEEEtran.cls version 1.8b or later) with an IEEE
%% journal paper.
%%
%% Support sites:
%% http://www.michaelshell.org/tex/ieeetran/
%% http://www.ctan.org/pkg/ieeetran
%% and
%% http://www.ieee.org/

%%*************************************************************************
%% Legal Notice:
%% This code is offered as-is without any warranty either expressed or
%% implied; without even the implied warranty of MERCHANTABILITY or
%% FITNESS FOR A PARTICULAR PURPOSE! 
%% User assumes all risk.
%% In no event shall the IEEE or any contributor to this code be liable for
%% any damages or losses, including, but not limited to, incidental,
%% consequential, or any other damages, resulting from the use or misuse
%% of any information contained here.
%%
%% All comments are the opinions of their respective authors and are not
%% necessarily endorsed by the IEEE.
%%
%% This work is distributed under the LaTeX Project Public License (LPPL)
%% ( http://www.latex-project.org/ ) version 1.3, and may be freely used,
%% distributed and modified. A copy of the LPPL, version 1.3, is included
%% in the base LaTeX documentation of all distributions of LaTeX released
%% 2003/12/01 or later.
%% Retain all contribution notices and credits.
%% ** Modified files should be clearly indicated as such, including  **
%% ** renaming them and changing author support contact information. **
%%*************************************************************************


% *** Authors should verify (and, if needed, correct) their LaTeX system  ***
% *** with the testflow diagnostic prior to trusting their LaTeX platform ***
% *** with production work. The IEEE's font choices and paper sizes can   ***
% *** trigger bugs that do not appear when using other class files.       ***                          ***
% The testflow support page is at:
% http://www.michaelshell.org/tex/testflow/



\documentclass[journal]{IEEEtran}
%
% If IEEEtran.cls has not been installed into the LaTeX system files,
% manually specify the path to it like:
% \documentclass[journal]{../sty/IEEEtran}





% Some very useful LaTeX packages include:
% (uncomment the ones you want to load)


% *** MISC UTILITY PACKAGES ***
%
%\usepackage{ifpdf}
% Heiko Oberdiek's ifpdf.sty is very useful if you need conditional
% compilation based on whether the output is pdf or dvi.
% usage:
% \ifpdf
%   % pdf code
% \else
%   % dvi code
% \fi
% The latest version of ifpdf.sty can be obtained from:
% http://www.ctan.org/pkg/ifpdf
% Also, note that IEEEtran.cls V1.7 and later provides a builtin
% \ifCLASSINFOpdf conditional that works the same way.
% When switching from latex to pdflatex and vice-versa, the compiler may
% have to be run twice to clear warning/error messages.






% *** CITATION PACKAGES ***
%
%\usepackage{cite}
% cite.sty was written by Donald Arseneau
% V1.6 and later of IEEEtran pre-defines the format of the cite.sty package
% \cite{} output to follow that of the IEEE. Loading the cite package will
% result in citation numbers being automatically sorted and properly
% "compressed/ranged". e.g., [1], [9], [2], [7], [5], [6] without using
% cite.sty will become [1], [2], [5]--[7], [9] using cite.sty. cite.sty's
% \cite will automatically add leading space, if needed. Use cite.sty's
% noadjust option (cite.sty V3.8 and later) if you want to turn this off
% such as if a citation ever needs to be enclosed in parenthesis.
% cite.sty is already installed on most LaTeX systems. Be sure and use
% version 5.0 (2009-03-20) and later if using hyperref.sty.
% The latest version can be obtained at:
% http://www.ctan.org/pkg/cite
% The documentation is contained in the cite.sty file itself.






% *** GRAPHICS RELATED PACKAGES ***
%
\ifCLASSINFOpdf
  % \usepackage[pdftex]{graphicx}
  % declare the path(s) where your graphic files are
  % \graphicspath{{../pdf/}{../jpeg/}}
  % and their extensions so you won't have to specify these with
  % every instance of \includegraphics
  % \DeclareGraphicsExtensions{.pdf,.jpeg,.png}
\else
  % or other class option (dvipsone, dvipdf, if not using dvips). graphicx
  % will default to the driver specified in the system graphics.cfg if no
  % driver is specified.
  % \usepackage[dvips]{graphicx}
  % declare the path(s) where your graphic files are
  % \graphicspath{{../eps/}}
  % and their extensions so you won't have to specify these with
  % every instance of \includegraphics
  % \DeclareGraphicsExtensions{.eps}
\fi
% graphicx was written by David Carlisle and Sebastian Rahtz. It is
% required if you want graphics, photos, etc. graphicx.sty is already
% installed on most LaTeX systems. The latest version and documentation
% can be obtained at: 
% http://www.ctan.org/pkg/graphicx
% Another good source of documentation is "Using Imported Graphics in
% LaTeX2e" by Keith Reckdahl which can be found at:
% http://www.ctan.org/pkg/epslatex
%
% latex, and pdflatex in dvi mode, support graphics in encapsulated
% postscript (.eps) format. pdflatex in pdf mode supports graphics
% in .pdf, .jpeg, .png and .mps (metapost) formats. Users should ensure
% that all non-photo figures use a vector format (.eps, .pdf, .mps) and
% not a bitmapped formats (.jpeg, .png). The IEEE frowns on bitmapped formats
% which can result in "jaggedy"/blurry rendering of lines and letters as
% well as large increases in file sizes.
%
% You can find documentation about the pdfTeX application at:
% http://www.tug.org/applications/pdftex





% *** MATH PACKAGES ***
%
%\usepackage{amsmath}
% A popular package from the American Mathematical Society that provides
% many useful and powerful commands for dealing with mathematics.
%
% Note that the amsmath package sets \interdisplaylinepenalty to 10000
% thus preventing page breaks from occurring within multiline equations. Use:
%\interdisplaylinepenalty=2500
% after loading amsmath to restore such page breaks as IEEEtran.cls normally
% does. amsmath.sty is already installed on most LaTeX systems. The latest
% version and documentation can be obtained at:
% http://www.ctan.org/pkg/amsmath





% *** SPECIALIZED LIST PACKAGES ***
%
%\usepackage{algorithmic}
% algorithmic.sty was written by Peter Williams and Rogerio Brito.
% This package provides an algorithmic environment fo describing algorithms.
% You can use the algorithmic environment in-text or within a figure
% environment to provide for a floating algorithm. Do NOT use the algorithm
% floating environment provided by algorithm.sty (by the same authors) or
% algorithm2e.sty (by Christophe Fiorio) as the IEEE does not use dedicated
% algorithm float types and packages that provide these will not provide
% correct IEEE style captions. The latest version and documentation of
% algorithmic.sty can be obtained at:
% http://www.ctan.org/pkg/algorithms
% Also of interest may be the (relatively newer and more customizable)
% algorithmicx.sty package by Szasz Janos:
% http://www.ctan.org/pkg/algorithmicx




% *** ALIGNMENT PACKAGES ***
%
%\usepackage{array}
% Frank Mittelbach's and David Carlisle's array.sty patches and improves
% the standard LaTeX2e array and tabular environments to provide better
% appearance and additional user controls. As the default LaTeX2e table
% generation code is lacking to the point of almost being broken with
% respect to the quality of the end results, all users are strongly
% advised to use an enhanced (at the very least that provided by array.sty)
% set of table tools. array.sty is already installed on most systems. The
% latest version and documentation can be obtained at:
% http://www.ctan.org/pkg/array


% IEEEtran contains the IEEEeqnarray family of commands that can be used to
% generate multiline equations as well as matrices, tables, etc., of high
% quality.




% *** SUBFIGURE PACKAGES ***
%\ifCLASSOPTIONcompsoc
%  \usepackage[caption=false,font=normalsize,labelfont=sf,textfont=sf]{subfig}
%\else
%  \usepackage[caption=false,font=footnotesize]{subfig}
%\fi
% subfig.sty, written by Steven Douglas Cochran, is the modern replacement
% for subfigure.sty, the latter of which is no longer maintained and is
% incompatible with some LaTeX packages including fixltx2e. However,
% subfig.sty requires and automatically loads Axel Sommerfeldt's caption.sty
% which will override IEEEtran.cls' handling of captions and this will result
% in non-IEEE style figure/table captions. To prevent this problem, be sure
% and invoke subfig.sty's "caption=false" package option (available since
% subfig.sty version 1.3, 2005/06/28) as this is will preserve IEEEtran.cls
% handling of captions.
% Note that the Computer Society format requires a larger sans serif font
% than the serif footnote size font used in traditional IEEE formatting
% and thus the need to invoke different subfig.sty package options depending
% on whether compsoc mode has been enabled.
%
% The latest version and documentation of subfig.sty can be obtained at:
% http://www.ctan.org/pkg/subfig




% *** FLOAT PACKAGES ***
%
%\usepackage{fixltx2e}
% fixltx2e, the successor to the earlier fix2col.sty, was written by
% Frank Mittelbach and David Carlisle. This package corrects a few problems
% in the LaTeX2e kernel, the most notable of which is that in current
% LaTeX2e releases, the ordering of single and double column floats is not
% guaranteed to be preserved. Thus, an unpatched LaTeX2e can allow a
% single column figure to be placed prior to an earlier double column
% figure.
% Be aware that LaTeX2e kernels dated 2015 and later have fixltx2e.sty's
% corrections already built into the system in which case a warning will
% be issued if an attempt is made to load fixltx2e.sty as it is no longer
% needed.
% The latest version and documentation can be found at:
% http://www.ctan.org/pkg/fixltx2e


%\usepackage{stfloats}
% stfloats.sty was written by Sigitas Tolusis. This package gives LaTeX2e
% the ability to do double column floats at the bottom of the page as well
% as the top. (e.g., "\begin{figure*}[!b]" is not normally possible in
% LaTeX2e). It also provides a command:
%\fnbelowfloat
% to enable the placement of footnotes below bottom floats (the standard
% LaTeX2e kernel puts them above bottom floats). This is an invasive package
% which rewrites many portions of the LaTeX2e float routines. It may not work
% with other packages that modify the LaTeX2e float routines. The latest
% version and documentation can be obtained at:
% http://www.ctan.org/pkg/stfloats
% Do not use the stfloats baselinefloat ability as the IEEE does not allow
% \baselineskip to stretch. Authors submitting work to the IEEE should note
% that the IEEE rarely uses double column equations and that authors should try
% to avoid such use. Do not be tempted to use the cuted.sty or midfloat.sty
% packages (also by Sigitas Tolusis) as the IEEE does not format its papers in
% such ways.
% Do not attempt to use stfloats with fixltx2e as they are incompatible.
% Instead, use Morten Hogholm'a dblfloatfix which combines the features
% of both fixltx2e and stfloats:
%
% \usepackage{dblfloatfix}
% The latest version can be found at:
% http://www.ctan.org/pkg/dblfloatfix




%\ifCLASSOPTIONcaptionsoff
%  \usepackage[nomarkers]{endfloat}
% \let\MYoriglatexcaption\caption
% \renewcommand{\caption}[2][\relax]{\MYoriglatexcaption[#2]{#2}}
%\fi
% endfloat.sty was written by James Darrell McCauley, Jeff Goldberg and 
% Axel Sommerfeldt. This package may be useful when used in conjunction with 
% IEEEtran.cls'  captionsoff option. Some IEEE journals/societies require that
% submissions have lists of figures/tables at the end of the paper and that
% figures/tables without any captions are placed on a page by themselves at
% the end of the document. If needed, the draftcls IEEEtran class option or
% \CLASSINPUTbaselinestretch interface can be used to increase the line
% spacing as well. Be sure and use the nomarkers option of endfloat to
% prevent endfloat from "marking" where the figures would have been placed
% in the text. The two hack lines of code above are a slight modification of
% that suggested by in the endfloat docs (section 8.4.1) to ensure that
% the full captions always appear in the list of figures/tables - even if
% the user used the short optional argument of \caption[]{}.
% IEEE papers do not typically make use of \caption[]'s optional argument,
% so this should not be an issue. A similar trick can be used to disable
% captions of packages such as subfig.sty that lack options to turn off
% the subcaptions:
% For subfig.sty:
% \let\MYorigsubfloat\subfloat
% \renewcommand{\subfloat}[2][\relax]{\MYorigsubfloat[]{#2}}
% However, the above trick will not work if both optional arguments of
% the \subfloat command are used. Furthermore, there needs to be a
% description of each subfigure *somewhere* and endfloat does not add
% subfigure captions to its list of figures. Thus, the best approach is to
% avoid the use of subfigure captions (many IEEE journals avoid them anyway)
% and instead reference/explain all the subfigures within the main caption.
% The latest version of endfloat.sty and its documentation can obtained at:
% http://www.ctan.org/pkg/endfloat
%
% The IEEEtran \ifCLASSOPTIONcaptionsoff conditional can also be used
% later in the document, say, to conditionally put the References on a 
% page by themselves.




% *** PDF, URL AND HYPERLINK PACKAGES ***
%
%\usepackage{url}
% url.sty was written by Donald Arseneau. It provides better support for
% handling and breaking URLs. url.sty is already installed on most LaTeX
% systems. The latest version and documentation can be obtained at:
% http://www.ctan.org/pkg/url
% Basically, \url{my_url_here}.




% *** Do not adjust lengths that control margins, column widths, etc. ***
% *** Do not use packages that alter fonts (such as pslatex).         ***
% There should be no need to do such things with IEEEtran.cls V1.6 and later.
% (Unless specifically asked to do so by the journal or conference you plan
% to submit to, of course. )


% correct bad hyphenation here
\hyphenation{op-tical net-works semi-conduc-tor}


\begin{document}
%
% paper title
% Titles are generally capitalized except for words such as a, an, and, as,
% at, but, by, for, in, nor, of, on, or, the, to and up, which are usually
% not capitalized unless they are the first or last word of the title.
% Linebreaks \\ can be used within to get better formatting as desired.
% Do not put math or special symbols in the title.
\title{A Review of Semantic Segmentation Models:\\ From Fully Convolutional Networks to\\ Promptable Foundation Models}
%
%
% author names and IEEE memberships
% note positions of commas and nonbreaking spaces ( ~ ) LaTeX will not break
% a structure at a ~ so this keeps an author's name from being broken across
% two lines.
% use \thanks{} to gain access to the first footnote area
% a separate \thanks must be used for each paragraph as LaTeX2e's \thanks
% was not built to handle multiple paragraphs
%

\author{Oscar~A.~Mart\'{i}nez~S.% <-this % stops a space
\thanks{Author affiliation goes here. E-mail contact to be completed by the author.}% <-this % stops a space
\thanks{Manuscript prepared October 15, 2026.}}

% note the % following the last \IEEEmembership and also \thanks - 
% these prevent an unwanted space from occurring between the last author name
% and the end of the author line. i.e., if you had this:
% 
% \author{....lastname \thanks{...} \thanks{...} }
%                     ^------------^------------^----Do not want these spaces!
%
% a space would be appended to the last name and could cause every name on that
% line to be shifted left slightly. This is one of those "LaTeX things". For
% instance, "\textbf{A} \textbf{B}" will typeset as "A B" not "AB". To get
% "AB" then you have to do: "\textbf{A}\textbf{B}"
% \thanks is no different in this regard, so shield the last } of each \thanks
% that ends a line with a % and do not let a space in before the next \thanks.
% Spaces after \IEEEmembership other than the last one are OK (and needed) as
% you are supposed to have spaces between the names. For what it is worth,
% this is a minor point as most people would not even notice if the said evil
% space somehow managed to creep in.



% The paper headers
\markboth{Journal of \LaTeX\ Class Files,~Vol.~14, No.~8, August~2015}%
{Mart\'{i}nez: A Review of Semantic Segmentation Models}
% The only time the second header will appear is for the odd numbered pages
% after the title page when using the twoside option.
% 
% *** Note that you probably will NOT want to include the author's ***
% *** name in the headers of peer review papers.                   ***
% You can use \ifCLASSOPTIONpeerreview for conditional compilation here if
% you desire.




% If you want to put a publisher's ID mark on the page you can do it like
% this:
%\IEEEpubid{0000--0000/00\$00.00~\copyright~2015 IEEE}
% Remember, if you use this you must call \IEEEpubidadjcol in the second
% column for its text to clear the IEEEpubid mark.



% use for special paper notices
%\IEEEspecialpapernotice{(Invited Paper)}




% make the title area
\maketitle

% As a general rule, do not put math, special symbols or citations
% in the abstract or keywords.
\begin{abstract}
Semantic segmentation, the assignment of a category label to every pixel of an
image, has progressed in under a decade from hand-tuned convolutional pipelines
to promptable foundation models trained on billions of masks. This paper reviews
five influential and highly cited segmentation models that mark the successive
turning points of that trajectory: the fully convolutional network, DeepLabv3+,
SegFormer, Mask2Former and the Segment Anything Model. For each model we identify
the architectural commitment it introduced, the evidence its authors reported, and
the assumptions that evidence rests on. Rather than tabulating benchmark scores in
isolation, we organise the literature into four thematic clusters---receptive-field
engineering, hierarchical attention, mask classification, and promptable
pre-training---and trace the methodological shift from per-pixel classification
towards set prediction and task-conditioned universal architectures. We find broad
consensus that mask classification subsumes per-pixel classification and that
large-scale pre-training dominates architectural refinement, but persistent
disagreement over how accuracy should be weighed against inference cost, and over
whether headline numbers obtained under differing backbones, crop sizes and
multi-scale inference protocols are comparable at all. We close by identifying open
questions concerning evaluation hygiene, boundary quality, label efficiency and
deployment under constrained compute.
\end{abstract}

% Note that keywords are not normally used for peerreview papers.
\begin{IEEEkeywords}
Semantic segmentation, scene understanding, convolutional neural networks, vision
transformers, mask classification, foundation models, literature review.
\end{IEEEkeywords}






% For peer review papers, you can put extra information on the cover
% page as needed:
% \ifCLASSOPTIONpeerreview
% \begin{center} \bfseries EDICS Category: 3-BBND \end{center}
% \fi
%
% For peerreview papers, this IEEEtran command inserts a page break and
% creates the second title. It will be ignored for other modes.
\IEEEpeerreviewmaketitle



\section{Introduction}
% The very first letter is a 2 line initial drop letter followed
% by the rest of the first word in caps.
% 
% form to use if the first word consists of a single letter:
% \IEEEPARstart{A}{demo} file is ....
% 
% form to use if you need the single drop letter followed by
% normal text (unknown if ever used by the IEEE):
% \IEEEPARstart{A}{}demo file is ....
% 
% Some journals put the first two words in caps:
% \IEEEPARstart{T}{his demo} file is ....
% 
% Here we have the typical use of a "T" for an initial drop letter
% and "HIS" in caps to complete the first word.
\IEEEPARstart{S}{emantic} segmentation assigns a category label to every pixel in
an image, producing a dense, structured description of a scene rather than a single
image-level decision. It underpins autonomous driving, medical image analysis,
remote sensing, robotic manipulation and photo editing, and it is the canonical
dense-prediction task against which advances in visual representation learning are
measured. Because the output lives at the resolution of the input, segmentation
models are simultaneously constrained by two competing demands: they must capture
long-range semantic context in order to decide \emph{what} a region is, and they
must preserve fine spatial detail in order to decide precisely \emph{where} it
begins and ends. Almost every architectural idea in the field can be read as a
particular compromise between those two demands.

The modern era of the task begins with the fully convolutional network (FCN) of
Long \emph{et al.}~\cite{long2014fully}, which reframed segmentation as per-pixel
classification and showed that a classification backbone could be trained
end-to-end, pixels-to-pixels, by replacing its fully connected layers with
convolutions. Subsequent work largely accepted that framing and attacked its
weaknesses. DeepLabv3+~\cite{chen2018encoderdecoder} addressed the loss of spatial
detail caused by repeated striding, combining atrous spatial pyramid pooling with a
decoder and depthwise separable convolutions. SegFormer~\cite{xie2021segformer}
replaced the convolutional backbone with a hierarchical, positional-encoding-free
transformer and a deliberately minimal all-MLP decoder. Mask2Former
\cite{cheng2021maskedattention} abandoned per-pixel classification altogether in
favour of mask classification, predicting a set of binary masks with associated
labels and thereby unifying semantic, instance and panoptic segmentation under one
architecture. OneFormer~\cite{jain2022oneformer} pushed that unification further
still, training a single model once and conditioning it on a task token at
inference. Finally, the Segment Anything Model (SAM)~\cite{kirillov2023segment}
changed the object of study: instead of training a model for a fixed label set, it
pre-trains a promptable model on over one billion masks and transfers zero-shot to
new distributions.

This paper reviews that trajectory through the lens of five models that define its
turning points, together with the FCN baseline from which they all descend. Our aim
is synthesis rather than enumeration. We ask what architectural commitment each
model made, what evidence its authors offered, what that evidence assumes, and
where the reported results genuinely agree or only appear to. The remainder of the
paper is organised as follows. Section~\ref{sec:sota} presents the state of the
art: it states the scope and inclusion criteria of the review, groups prior work
into thematic clusters, traces methodological trends over time, separates consensus
findings from unresolved disagreements, and identifies gaps and future directions.
Section~\ref{sec:conclusion} concludes.

\section{State of the Art}
\label{sec:sota}

\subsection{Scope and Inclusion Criteria}

This review covers learning-based models for dense semantic labelling of natural
RGB images. A work was included if it satisfied three criteria. First, it had to
introduce a general-purpose segmentation \emph{architecture or pre-training
paradigm}, not an application of an existing one to a particular domain; this
excludes the large body of work applying SegFormer or SAM to agriculture, remote
sensing, histopathology or industrial inspection, valuable though it is. Second, it
had to report results on at least one of the three standard public benchmarks ---
PASCAL VOC 2012, Cityscapes or ADE20K --- so that claims could be placed on a
common axis. Third, it had to have demonstrably shifted subsequent practice, judged
by whether later architectures in the set treat it as the baseline to beat.

Applying these criteria yields six works, listed in Table~\ref{tab:models}: the FCN
as a historical anchor, and five models that we treat as the current reference
points of the field. Three boundaries of this scope should be stated plainly.
Instance and panoptic segmentation enter the discussion only where a model claims
to unify them with semantic segmentation, as Mask2Former and OneFormer do. Video,
3D and medical-specific segmentation are out of scope. And because the field
publishes predominantly through arXiv preprints and conference proceedings, the
evidence reviewed here is drawn from author-reported results, which have not been
independently re-run for this review --- a limitation we return to in
Section~\ref{subsec:disagree}.

\begin{table*}[!t]
\renewcommand{\arraystretch}{1.3}
\caption{The Five Reviewed Models and Their Historical Anchor. Results Are as Reported by the Original Authors; See Section~\ref{subsec:disagree} on Comparability.}
\label{tab:models}
\centering
\begin{tabular}{|l|c|l|l|}
\hline
\textbf{Model} & \textbf{Year} & \textbf{Defining commitment} & \textbf{Headline result reported by authors}\\
\hline
FCN~\cite{long2014fully} & 2014 & Per-pixel classification, trained end-to-end & First end-to-end pixels-to-pixels segmentation network\\
\hline
DeepLabv3+~\cite{chen2018encoderdecoder} & 2018 & Atrous pyramid pooling plus a decoder & 89.0\% on PASCAL VOC 2012 test; 82.1\% on Cityscapes test\\
\hline
SegFormer~\cite{xie2021segformer} & 2021 & Hierarchical transformer, all-MLP decoder & B5: 84.0\% mIoU Cityscapes val; B0: 71.9\% mIoU at 48 FPS\\
\hline
Mask2Former~\cite{cheng2021maskedattention} & 2021 & Mask classification with masked attention & 57.7 mIoU ADE20K; 57.8 PQ and 50.1 AP on COCO\\
\hline
OneFormer~\cite{jain2022oneformer} & 2022 & Task-conditioned joint training, one model & Beats per-task Mask2Former using $\sim$3$\times$ fewer training resources\\
\hline
SAM~\cite{kirillov2023segment} & 2023 & Promptable pre-training on 1B masks & Zero-shot transfer often competitive with supervised baselines\\
\hline
\end{tabular}
\end{table*}

\subsection{Thematic Clusters in Prior Work}

The six works fall into four clusters, each defined by the mechanism it uses to
reconcile semantic context with spatial precision.

\subsubsection{Receptive-Field Engineering}
The first cluster accepts the FCN's per-pixel formulation and enlarges the
effective receptive field without destroying resolution. DeepLabv3+ is its mature
representative. Its authors frame the problem explicitly as a choice between two
families: spatial pyramid pooling, which encodes multi-scale context by probing
features at several dilation rates, and encoder--decoder structures, which recover
sharp boundaries by gradually restoring spatial information
\cite{chen2018encoderdecoder}. Their contribution is to refuse the choice and
combine both, using DeepLabv3 as the encoder and adding a simple decoder. The paper
also makes an argument that recurs throughout the field: extracting dense features
purely by dilation is computationally prohibitive, since obtaining an output stride
of 8 with ResNet-101 requires dilating 78 layers. Atrous separable convolution is
introduced precisely to make this affordable.

\subsubsection{Hierarchical Attention}
The second cluster replaces the convolutional backbone with a transformer while
retaining per-pixel classification. SegFormer is the clearest statement of this
position, and it is as much a critique of earlier vision transformers as an
architecture. Its authors identify two concrete defects in the ViT-based SETR
approach: ViT emits single-scale, low-resolution features, and its cost grows
sharply with image size \cite{xie2021segformer}. SegFormer answers with a
hierarchical encoder that emits multi-scale features and, notably, drops positional
encoding entirely --- a choice motivated not by elegance but by a specific failure
mode, namely the accuracy loss incurred when positional codes are interpolated for
test resolutions that differ from training. The second half of the argument is
about the decoder. Observing that attention in lower layers stays local while
attention in the highest layers is strongly non-local, the authors argue that
simply aggregating across layers with an MLP already combines local and global
evidence, making heavy decoders unnecessary.

\subsubsection{Mask Classification and Universal Architectures}
The third cluster rejects per-pixel classification. Mask2Former's framing is that
semantic, instance and panoptic segmentation differ only in what semantics are used
to group pixels, so specialising an architecture per task duplicates research and
hardware effort \cite{cheng2021maskedattention}. The technical core is masked
attention, which constrains cross-attention to the regions of the currently
predicted mask instead of attending to the whole image; the authors report that
this yields both faster convergence and higher accuracy. Three further changes are
reported as necessary rather than incidental: multi-scale high-resolution features
to recover small objects, a reordering of self- and cross-attention with learnable
queries, and computing the mask loss on a few randomly sampled points, which the
authors state saves $3\times$ training memory with no loss of performance.

This last point deserves emphasis because it is an accessibility claim, not only an
efficiency one. Mask2Former's authors note that its predecessor MaskFormer required
300 epochs to reach 40.1 AP and could fit only a single image on a 32\,GB GPU,
whereas the specialised Swin-HTC++ reached better accuracy in 72 epochs. The
argument is that universal architectures had been held back as much by training
cost as by accuracy.

OneFormer extends the cluster by attacking a residual inconsistency. Mask2Former is
universal in architecture but not in training: to obtain its best numbers it is
trained separately on each task, which on ADE20K means three runs of 160k
iterations and three sets of weights to store and serve
\cite{jain2022oneformer}. OneFormer's authors hypothesise that joint training fails
without explicit task guidance, and introduce a task token of the form ``the task is
\{task\}'' to condition the model, together with a query--text contrastive loss to
sharpen inter-task and inter-class distinctions. Because panoptic annotations
contain both semantic and instance information, semantic and instance labels are
derived from panoptic ground truth, so only panoptic data is needed at training
time.

\subsubsection{Promptable Pre-Training}
The fourth cluster changes the task definition rather than the architecture. SAM
proposes promptable segmentation, in which the model must return a valid mask for
any prompt --- a point, a box, a coarse mask or text \cite{kirillov2023segment}.
Ambiguity is handled by design: a point on a shirt might indicate the shirt or the
person, so the model predicts several masks and is required only to be right about
one. Architecturally the design is deliberately plain --- a heavy image encoder, a
light prompt encoder, and a mask decoder that runs in roughly 50\,ms in a browser
once the image embedding is cached. The substance of the contribution is the data
engine: because no web-scale mask corpus existed, the authors iterated between
using the model to assist annotation and using the resulting annotations to improve
the model, producing SA-1B with over one billion masks on 11 million images.

\subsection{Methodological Trends Over Time}

Read chronologically, four trends are visible across these clusters.

The first is the migration of the inductive bias from the operator to the data.
DeepLabv3+ encodes multi-scale reasoning structurally, in the dilation rates of its
pyramid. SegFormer moves that burden into the attention mechanism. SAM removes the
architectural commitment almost entirely and relies on the scale and diversity of
SA-1B. The architectures do not become more sophisticated over time; if anything
they become simpler, and the complexity migrates into training data and objectives.

The second is the displacement of per-pixel classification by set prediction. The
FCN formulation held for roughly seven years. Its replacement by mask
classification was decisive once Mask2Former demonstrated that a single mask-based
architecture could match or beat specialised models on all three tasks.

The third is the steady relocation of the unit of unification: first across tasks
with one architecture (Mask2Former), then across tasks with one trained model
(OneFormer), then across data distributions with one pre-trained model (SAM). Each
step widens the claim of generality and, correspondingly, weakens the guarantee of
optimality on any particular benchmark.

The fourth is the rise of efficiency as a reported quantity rather than an
afterthought. SegFormer reports parameters, FLOPs and frame rate alongside
accuracy, pointing out that B5 is $4\times$ smaller than SETR while more accurate,
and that B0 reaches 71.9\% mIoU at 48 FPS on Cityscapes. OneFormer frames its
contribution partly in terms of a $3\times$ reduction in training time and storage.
This reflects a field increasingly conscious that benchmark gains purchased with
unbounded compute are not deployable gains.

\subsection{Consensus Findings and Unresolved Disagreements}
\label{subsec:disagree}

Several findings are now well established across independent groups. Mask
classification generalises per-pixel classification: a single mask-based
architecture matches or exceeds specialised per-task architectures on semantic,
instance and panoptic segmentation \cite{cheng2021maskedattention,jain2022oneformer}.
Multi-scale features are necessary for dense prediction regardless of whether the
backbone is convolutional or attention-based; every model reviewed here either
builds them in or is criticised by its successors for failing to
\cite{chen2018encoderdecoder,xie2021segformer}. Decoder complexity is largely
unnecessary once the encoder produces multi-scale features, a conclusion SegFormer
argues directly and Mask2Former corroborates from the opposite direction. And
large-scale pre-training transfers: SAM's zero-shot performance is reported as
frequently competitive with, and sometimes superior to, fully supervised results
\cite{kirillov2023segment}.

Three disagreements remain genuinely open, and we state them separately from the
findings above because the evidence does not settle them.

The first is comparability. Headline numbers in this literature are reported under
different backbones, crop sizes, iteration budgets, query counts and inference
protocols. OneFormer's own comparison tables make this visible: entries differ in
parameters, FLOPs, crop size and training iterations, and report single-scale and
multi-scale mIoU in separate columns \cite{jain2022oneformer}. Mask2Former's
reported 57.7 mIoU on ADE20K and the 54.5 mIoU attributed to a Mask2Former
configuration in OneFormer's table are not in conflict --- they describe different
settings --- but a reader comparing abstracts alone would conclude otherwise. The
field has no agreed protocol for reporting these axes together, and claims of
state-of-the-art status should be treated as conditional on the configuration until
verified.

The second concerns the accuracy--efficiency frontier. There is no consensus on how
to weigh the two. SegFormer's position is that a model is only state-of-the-art if
it is simultaneously efficient; Mask2Former and OneFormer optimise primarily for
accuracy and treat training cost, not inference cost, as the efficiency axis worth
reporting. These are different objectives presented in the same units, and the
literature does not adjudicate between them.

The third concerns the limits of foundation models. SAM's authors are explicit that
the model is not uniformly superior: it can miss fine structures, sometimes
hallucinates small disconnected components, and produces less crisp boundaries than
more computationally intensive methods that zoom in. They further expect dedicated
interactive segmentation methods to outperform SAM when many prompt points are
available, noting that SAM is designed for generality and breadth rather than high
interactive IoU \cite{kirillov2023segment}. Whether promptable pre-training will
eventually dominate task-specific training on accuracy, or settle into a
complementary role as an initialisation and annotation tool, is not yet answerable
from the published evidence.

\subsection{Gaps and Open Research Questions}

Four gaps follow from the preceding analysis.

\emph{Evaluation hygiene.} The comparability problem is methodological, not
technical, and it is the most tractable. The field lacks a standard reporting
convention that pins backbone, pre-training corpus, crop size, schedule length and
single- versus multi-scale inference to every reported score. Without it, measured
progress and measured compute cannot be separated.

\emph{Boundary quality.} Every model reviewed here reports mIoU, a region-overlap
measure that is relatively insensitive to boundary error. Yet boundary fidelity is
exactly what DeepLabv3+'s decoder was designed to improve, and exactly where SAM's
authors concede weakness. A metric the community agrees on would make an
acknowledged weakness measurable.

\emph{Label efficiency.} SAM demonstrates that a data engine can produce masks at
scale, but SA-1B's masks are class-agnostic. Semantic segmentation needs category
labels, and how to obtain them at comparable scale without proportional annotation
cost is unresolved.

\emph{Deployment under constrained compute.} The accuracy frontier is occupied by
models whose cost excludes most embedded and real-time applications. SegFormer-B0
shows the frontier can be approached cheaply, but no systematic account exists of
what is lost in doing so.

\subsection{Suggested Future Directions}

We suggest four directions, ordered by how directly the reviewed evidence supports
them.

First, a standardised evaluation protocol and a shared reporting template. This
requires no new modelling and would resolve the comparability disagreement
described above; it is the highest-value, lowest-risk contribution available.

Second, distillation of promptable foundation models into task-specific, deployable
students. SAM's design already separates an expensive image encoder from a
lightweight decoder, which makes the encoder the natural target for compression;
early work in this direction is appearing, and the architectural separation
suggests it should succeed.

Third, closing the loop between promptable pre-training and universal
architectures. SAM supplies class-agnostic masks at scale; OneFormer supplies
task-conditioned semantic reasoning. Their combination --- using SAM-scale
pre-training to initialise a task-conditioned universal model --- is suggested by
the complementarity of the two, though we note this is a tentative inference from
the design of the two systems rather than a conclusion supported by published
results.

Fourth, boundary-aware objectives and metrics. Given that boundary error is both an
acknowledged weakness and an unmeasured one, progress here would be visible
immediately once the metric exists.



% An example of a floating figure using the graphicx package.
% Note that \label must occur AFTER (or within) \caption.
% For figures, \caption should occur after the \includegraphics.
% Note that IEEEtran v1.7 and later has special internal code that
% is designed to preserve the operation of \label within \caption
% even when the captionsoff option is in effect. However, because
% of issues like this, it may be the safest practice to put all your
% \label just after \caption rather than within \caption{}.
%
% Reminder: the "draftcls" or "draftclsnofoot", not "draft", class
% option should be used if it is desired that the figures are to be
% displayed while in draft mode.
%
%\begin{figure}[!t]
%\centering
%\includegraphics[width=2.5in]{myfigure}
% where an .eps filename suffix will be assumed under latex, 
% and a .pdf suffix will be assumed for pdflatex; or what has been declared
% via \DeclareGraphicsExtensions.
%\caption{Simulation results for the network.}
%\label{fig_sim}
%\end{figure}

% Note that the IEEE typically puts floats only at the top, even when this
% results in a large percentage of a column being occupied by floats.


% An example of a double column floating figure using two subfigures.
% (The subfig.sty package must be loaded for this to work.)
% The subfigure \label commands are set within each subfloat command,
% and the \label for the overall figure must come after \caption.
% \hfil is used as a separator to get equal spacing.
% Watch out that the combined width of all the subfigures on a 
% line do not exceed the text width or a line break will occur.
%
%\begin{figure*}[!t]
%\centering
%\subfloat[Case I]{\includegraphics[width=2.5in]{box}%
%\label{fig_first_case}}
%\hfil
%\subfloat[Case II]{\includegraphics[width=2.5in]{box}%
%\label{fig_second_case}}
%\caption{Simulation results for the network.}
%\label{fig_sim}
%\end{figure*}
%
% Note that often IEEE papers with subfigures do not employ subfigure
% captions (using the optional argument to \subfloat[]), but instead will
% reference/describe all of them (a), (b), etc., within the main caption.
% Be aware that for subfig.sty to generate the (a), (b), etc., subfigure
% labels, the optional argument to \subfloat must be present. If a
% subcaption is not desired, just leave its contents blank,
% e.g., \subfloat[].


% An example of a floating table. Note that, for IEEE style tables, the
% \caption command should come BEFORE the table and, given that table
% captions serve much like titles, are usually capitalized except for words
% such as a, an, and, as, at, but, by, for, in, nor, of, on, or, the, to
% and up, which are usually not capitalized unless they are the first or
% last word of the caption. Table text will default to \footnotesize as
% the IEEE normally uses this smaller font for tables.
% The \label must come after \caption as always.
%
%\begin{table}[!t]
%% increase table row spacing, adjust to taste
%\renewcommand{\arraystretch}{1.3}
% if using array.sty, it might be a good idea to tweak the value of
% \extrarowheight as needed to properly center the text within the cells
%\caption{An Example of a Table}
%\label{table_example}
%\centering
%% Some packages, such as MDW tools, offer better commands for making tables
%% than the plain LaTeX2e tabular which is used here.
%\begin{tabular}{|c||c|}
%\hline
%One & Two\\
%\hline
%Three & Four\\
%\hline
%\end{tabular}
%\end{table}


% Note that the IEEE does not put floats in the very first column
% - or typically anywhere on the first page for that matter. Also,
% in-text middle ("here") positioning is typically not used, but it
% is allowed and encouraged for Computer Society conferences (but
% not Computer Society journals). Most IEEE journals/conferences use
% top floats exclusively. 
% Note that, LaTeX2e, unlike IEEE journals/conferences, places
% footnotes above bottom floats. This can be corrected via the
% \fnbelowfloat command of the stfloats package.




\section{Conclusion}
\label{sec:conclusion}
Across the decade separating the fully convolutional network from the Segment
Anything Model, semantic segmentation has not so much refined one idea as replaced
its foundations twice: first when mask classification displaced per-pixel
classification, and again when large-scale promptable pre-training displaced
architectural specialisation as the main source of generalisation. The reviewed
evidence supports both shifts clearly. It supports far less clearly the headline
comparisons through which they are usually narrated, because reported scores depend
on backbone, schedule and inference protocol in ways that abstracts rarely disclose.
The most useful near-term contribution to this literature may therefore not be a new
architecture but a shared discipline for measuring the ones that exist.




% if have a single appendix:
%\appendix[Proof of the Zonklar Equations]
% or
%\appendix  % for no appendix heading
% do not use \section anymore after \appendix, only \section*
% is possibly needed

% use appendices with more than one appendix
% then use \section to start each appendix
% you must declare a \section before using any
% \subsection or using \label (\appendices by itself
% starts a section numbered zero.)
%


% use section* for acknowledgment
\section*{Acknowledgment}


The author thanks the maintainers of the arXiv open-access archive, without which
a review of this kind would not be practical.


% Can use something like this to put references on a page
% by themselves when using endfloat and the captionsoff option.
\ifCLASSOPTIONcaptionsoff
  \newpage
\fi



% trigger a \newpage just before the given reference
% number - used to balance the columns on the last page
% adjust value as needed - may need to be readjusted if
% the document is modified later
%\IEEEtriggeratref{8}
% The "triggered" command can be changed if desired:
%\IEEEtriggercmd{\enlargethispage{-5in}}

% references section

% can use a bibliography generated by BibTeX as a .bbl file
% BibTeX documentation can be easily obtained at:
% http://mirror.ctan.org/biblio/bibtex/contrib/doc/
% The IEEEtran BibTeX style support page is at:
% http://www.michaelshell.org/tex/ieeetran/bibtex/
%\bibliographystyle{IEEEtran}
% argument is your BibTeX string definitions and bibliography database(s)
%\bibliography{IEEEabrv,../bib/paper}
%
% <OR> manually copy in the resultant .bbl file
% set second argument of \begin to the number of references
% (used to reserve space for the reference number labels box)
\begin{thebibliography}{10}

\bibitem{long2014fully}
J.~Long, E.~Shelhamer, and T.~Darrell, ``Fully convolutional networks for
  semantic segmentation,'' \emph{arXiv preprint arXiv:1411.4038}, 2014.

\bibitem{chen2018encoderdecoder}
L.-C. Chen, Y.~Zhu, G.~Papandreou, F.~Schroff, and H.~Adam,
  ``Encoder-decoder with atrous separable convolution for semantic image
  segmentation,'' \emph{arXiv preprint arXiv:1802.02611}, 2018.

\bibitem{xie2021segformer}
E.~Xie, W.~Wang, Z.~Yu, A.~Anandkumar, J.~M. Alvarez, and P.~Luo,
  ``SegFormer: Simple and efficient design for semantic segmentation with
  transformers,'' \emph{arXiv preprint arXiv:2105.15203}, 2021.

\bibitem{cheng2021maskedattention}
B.~Cheng, I.~Misra, A.~G. Schwing, A.~Kirillov, and R.~Girdhar,
  ``Masked-attention mask transformer for universal image segmentation,''
  \emph{arXiv preprint arXiv:2112.01527}, 2021.

\bibitem{jain2022oneformer}
J.~Jain, J.~Li, M.~Chiu, A.~Hassani, N.~Orlov, and H.~Shi, ``OneFormer: One
  transformer to rule universal image segmentation,'' \emph{arXiv preprint
  arXiv:2211.06220}, 2022.

\bibitem{kirillov2023segment}
A.~Kirillov, E.~Mintun, N.~Ravi, H.~Mao, C.~Rolland, L.~Gustafson, T.~Xiao,
  S.~Whitehead, A.~C. Berg, W.-Y. Lo, P.~Doll\'{a}r, and R.~Girshick, ``Segment
  anything,'' \emph{arXiv preprint arXiv:2304.02643}, 2023.

\end{thebibliography}

% biography section
% 
% If you have an EPS/PDF photo (graphicx package needed) extra braces are
% needed around the contents of the optional argument to biography to prevent
% the LaTeX parser from getting confused when it sees the complicated
% \includegraphics command within an optional argument. (You could create
% your own custom macro containing the \includegraphics command to make things
% simpler here.)
%\begin{IEEEbiography}[{\includegraphics[width=1in,height=1.25in,clip,keepaspectratio]{mshell}}]{Michael Shell}
% or if you just want to reserve a space for a photo:

\begin{IEEEbiographynophoto}{Oscar A. Mart\'{i}nez S.}
Biography text here.
\end{IEEEbiographynophoto}

% You can push biographies down or up by placing
% a \vfill before or after them. The appropriate
% use of \vfill depends on what kind of text is
% on the last page and whether or not the columns
% are being equalized.

%\vfill

% Can be used to pull up biographies so that the bottom of the last one
% is flush with the other column.
%\enlargethispage{-5in}



% that's all folks
\end{document}
